%\usepackage[T1]{fontenc}
\usepackage{graphicx}
\usepackage{float}
\usepackage{verbatim}
\usepackage{enumerate}
\usepackage[colorlinks=true, linkcolor=blue]{hyperref}
\usepackage{amsmath, amsthm, amssymb, amscd}
\usepackage{polski}
\usepackage{caption}
\captionsetup[table]{justification=centering}
\captionsetup[figure]{justification=centering}
\newtheorem{thm}{Twierdzenie}
\newtheorem{lem}{Lemat}
\newtheorem{cor}{Wniosek}
\newtheorem{prop}{Propozycja}
\newtheorem{rem}{Uwaga}
\newtheorem{note}{Notka}
\newtheorem{fact}{Fakt}
\newtheorem{defn}{Definicja}
\renewcommand{\figurename}{Wykres}
\renewcommand{\tablename}{Tabela}
\counterwithin{equation}{section}
\def \E {\mathbb{E}}
\def \P {\mathbb{P}}

\usepackage[top=2.5cm, bottom=2.5cm, left=2cm, right=2cm]{geometry}
\makeatletter
\newcommand\numerraportu[1]{\renewcommand\@title{Raport #1}}
\renewcommand{\maketitle}{
	\begin{center}
		\vspace{5pt}
		\textbf{\Huge \sffamily Analiza szeregów czasowych}\\
		\vspace{15pt}
		\textbf{\Large \sffamily \@title}\\
		\vspace{10pt}
		{\Large \sffamily Autorzy: \@author}\\
		\vspace{10pt}
		{\large\sffamily \@date}\\
		\hrulefill\\
		\vspace{20pt}
	\end{center}
}
\usepackage{fancyhdr}
\pagestyle{fancy}
\fancyhf{}
\fancyhead[R]{\thepage}
\fancyhead[L]{\@author}
\setlength{\headheight}{15pt}
\setlength{\headsep}{15pt}
\makeatother

\thispagestyle{plain}